\section{Adaptive integration}\label{sec:adaptivity}

\pkg{} provides two complementary adaptive interfaces.  A generic step-doubling
controller supplies a reference mechanism applicable to one-step
\texttt{MethodSpec} objects.  A specialized solver implements the embedded
ETD34 pair of Whalen, Brio and Moloney~\cite{WhalenBrioMoloney2015}, built on
the scheme of Krogstad~\cite{Krogstad2005}, and adds backend-aware rejection
behaviour.  Their structural contracts are compared in
Table~\ref{tab:adaptivity}.

\input{table03_adaptivity.tex}

\subsection{Generic step doubling}

For a one-step method of order $p$, the generic controller compares a full step
of length $h$ with two steps of length $h/2$.  Its error indicator is
\begin{equation}
  e_n=\frac{y_{h/2,h/2}-y_h}{2^p-1}.
\end{equation}
The implementation accepts the two-half-step state; Richardson extrapolation is
not the default.  This route is deliberately simple and broad: it provides an
adaptive reference for any one-step catalogue entry without requiring a bespoke
embedded formula.  Its cost is correspondingly high.  A completed Krogstad
trial evaluates three complete method steps and twelve nonlinear values; the
first stage of the full step and of the first half step coincide, so eleven
would suffice, and the generic implementation does not exploit this.  The
step-doubling results of Section~\ref{sec:experiments} should be read with this
unoptimized baseline in mind.

Variable-step history methods are not silently routed through this controller.
The published interface excludes variable-step multistep and general linear
methods, because changing the step size also changes the representation and
transport of history data.  Fixed-step history methods remain fully supported.

\subsection{Embedded ETD34}

The specialized solver advances with the fourth-order Krogstad candidate, whose
stiff order is three in the sense of Hochbruck and
Ostermann~\cite{HochbruckOstermann2005}, and uses an embedded third-order
estimator.  With $N_4$ the nonlinear value at the last internal Krogstad stage
and $N_5=N(t_n+h,y_{n+1})$ evaluated at the fourth-order candidate, the
indicator is
\begin{equation}
  e_n=h\left[-\varphi_2(hL)+4\varphi_3(hL)\right](N_4-N_5),
  \label{eq:etd34-estimator}
\end{equation}
which is the same expression used by the \texttt{rkstiff} implementation of the
pair by its first author.  A completed trial requires four new nonlinear
evaluations after the initial endpoint value.  On acceptance, $N_5$ is reused
at the next step; on rejection, the old endpoint value is retained.  The scaled
error norm uses user-specified relative and absolute tolerances, a fourth-root
controller reflects the $O(h^4)$ local scaling of the estimate, and the
accepted state remains the fourth-order candidate.

The estimator has a precise limitation.  Both nonlinear values in
Eq.~\eqref{eq:etd34-estimator} are evaluated at $t_n+h$.  If
$N(t,y)=g(t)$ is independent of the state, their difference is zero even when
the step has nonzero error.  A self-contained example is $L=0$,
$y'=(1+t)^4$, $y(0)=0$.  Here the Krogstad step reduces to classical RK4,
so the first step and the exact solution satisfy
\begin{equation}
\begin{aligned}
y_1^{[4]}&=\frac h6\bigl[1+4(1+h/2)^4+(1+h)^4\bigr],\\
y(h)&=\frac{(1+h)^5-1}{5},\qquad
y_1^{[4]}-y(h)=\frac{h^5}{120},\qquad e_0=0.
\end{aligned}\label{eq:etd34-forcing-counterexample}
\end{equation}
This is an algebraic counterexample, not a new time-integration run.  It
exhibits the estimator degeneracy relevant to the robustness discussion of
Zoto and Bowman~\cite{ZotoBowman2023}.

For the more general form $N(t,y)=g(t)+F(y)$, only the direct additive value
$g(t_n+h)$ cancels.  The indicator still contains
$F(Y_4)-F(y_{n+1})$.  With $d=Y_4-y_{n+1}$ and $F$ differentiable along the
segment between the states,
\begin{equation}
 F(Y_4)-F(y_{n+1})=
 \left[\int_0^1 DF(y_{n+1}+\theta d)\,d\theta\right]d.
 \label{eq:forcing-indirect}
\end{equation}
The forcing affects the states and their difference and can therefore affect
the indicator indirectly.  Exact cancellation of the direct forcing term is
not a proof that the estimator is blind to all forcing-induced errors when
$F$ is nonconstant.  Conversely, state dependence alone does not establish
robust error detection.  The forced comparisons in
Section~\ref{sec:adaptivework} and the autonomous example in
Section~\ref{sec:broadband} are interpreted case by case.  Step doubling is the
documented alternative when the embedded indicator is unsuitable; replacing
the published pair is not a contribution of this article.

\subsection{Temporal and matrix-action tolerances}

An exponential integrator has two distinct accuracy mechanisms: temporal error
control and numerical approximation of matrix-function actions.  ETD34
therefore creates a work backend whose action tolerance is
\begin{equation}
 \tau_{\mathrm{act}}
 =\min(\tau_{\max},\max(\tau_{\min},f_{\mathrm{act}}\,\mathrm{rtol})),
 \label{eq:actiontol}
\end{equation}
where $f_{\mathrm{act}}$ is an explicit action-tolerance fraction.  The package
default is 0.05; the canonical timing campaign deliberately uses 0.01 and
records that value as part of the experiment configuration.  Neither number is
claimed to be universally optimal.  Supplementary Figure~S3 retains the
tolerance-coupling sweep that motivated this separation.  The work backend is
obtained from the user's backend through a single hook, \texttt{with\_tolerance},
so that external backends participate in the adaptive solvers on the same
footing as the built-in ones, and a calibrated selector keeps its calibration
record when cloned.

\subsection{Failed actions, non-finite trials and termination}

Iterative matrix-function algorithms can fail to meet their stopping criterion
for a particular trial step.  In ETD34, convergence exceptions from Leja,
KIOPS-type or Krylov actions, and trials that produce non-finite values, are
treated as rejected time steps: the solver halves the step and retries rather
than terminating the integration.  Both kinds of rejection are recorded
separately from ordinary error-estimator rejections, and the nonlinear
evaluations and matrix-vector products spent on failed attempts are counted,
not estimated.

Termination is explicit.  Both adaptive solvers use the minimum step
$\max(h_{\min},16\varepsilon T)$ with $T=\max(|t_0|,|t_f|,t_f-t_0)$, and an
attempt that fails or is rejected at that minimum ends the integration with an
error that names the cause.  A maximum step below the smallest resolvable step
of the interval is rejected before integrating, a step that would not advance
the floating-point time is reported, and a failure that leaves a step too
small to reach $t_f$ within the attempt budget stops the integration
immediately instead of exhausting the budget.  All option values are validated,
including finiteness.  These rules replace the behaviour of the first release,
in which a persistent failure at the minimum step consumed the whole attempt
budget.

\subsection{Cost-aware control is experimental}

The implementation also contains a cost-aware proposal based on observed work
per unit simulated time.  It can constrain the accuracy-based step proposal
when a local log--log cost gradient suggests a cheaper step size.  The work of
an attempt is the increment of the backend's own matrix-vector counter, and
the controller refuses to run in matvec units with a backend that reports no
counter; it never substitutes wall-clock time for a missing count.  The stress
test of Section~\ref{sec:experiments} shows why this feature is not a headline
performance claim: the matvec proxy decreases while recorded wall time can
increase, because the controller takes more steps and nonlinear evaluations.
We therefore retain it as an instrumented diagnostic, not as evidence of lower
solve time.
